% ===== PRÉAMBULE CANONIQUE — à coller VERBATIM en tête de CHAQUE .tex =====
\documentclass[11pt,a4paper]{article}
\usepackage[utf8]{inputenc}\usepackage[T1]{fontenc}\usepackage{lmodern}
\usepackage[french]{babel}
\usepackage{amsmath,amssymb,mathtools}\usepackage{icomma}
\usepackage[version=4]{mhchem}          % \ce{...} : équations & formules
\usepackage{siunitx}\sisetup{locale=FR} % \qty{}{}, \si{}, \ang{}
\usepackage{chemfig}                    % \chemfig{...} : structures développées
\usepackage{xcolor}\usepackage{colortbl}
\usepackage{tikz}\usepackage{pgfplots}\pgfplotsset{compat=1.18}
\usepackage[most]{tcolorbox}\usepackage{enumitem}\usepackage{array,booktabs}
\usepackage[a4paper,margin=2.2cm]{geometry}
\usetikzlibrary{arrows.meta,calc,angles,quotes,positioning,patterns,backgrounds,shapes.geometric}
\definecolor{primary}{RGB}{222,49,99}\definecolor{primarydark}{RGB}{150,22,55}
\definecolor{defcol}{RGB}{0,114,178}\definecolor{methcol}{RGB}{52,92,148}
\definecolor{excol}{RGB}{0,135,90}\definecolor{attncol}{RGB}{200,40,55}
\tcbset{boxbase/.style={enhanced,breakable,boxrule=0.9pt,arc=2.5pt,
  left=3mm,right=3mm,top=2.3mm,bottom=2.3mm,fonttitle=\bfseries,coltitle=white}}
% --- boîtes TUTORIEL (noms EXACTS, ne pas renommer) ---
\newtcolorbox{cle}[1][]{boxbase,colback=defcol!6,colframe=defcol,
  title={\textsc{À retenir}\ifx\relax#1\relax\relax\else\textmd{ : \itshape #1}\fi}}
\newtcolorbox{methode}[1][]{boxbase,colback=methcol!7,colframe=methcol,sharp corners,
  title={\textsc{Méthode}\ifx\relax#1\relax\relax\else\textmd{ --- \itshape #1}\fi}}
\newtcolorbox{exemplebox}[1][]{boxbase,colback=excol!5,colframe=excol,
  title={\textsc{Exemple}\ifx\relax#1\relax\relax\else\textmd{ : \itshape #1}\fi}}
\newtcolorbox{piege}[1][]{boxbase,colback=attncol!6,colframe=attncol,
  title={\textsc{Piège}\ifx\relax#1\relax\relax\else\textmd{ : \itshape #1}\fi}}
\newtcolorbox{remarque}[1][]{boxbase,colback=black!4,colframe=black!45,
  title={\textsc{Remarque}\ifx\relax#1\relax\relax\else\textmd{ : \itshape #1}\fi}}
% --- boîtes EXERCICE / CORRIGÉ (noms EXACTS, auto-comptées) ---
\newtcolorbox[auto counter]{exo}[1][]{boxbase,colback=primary!4,colframe=primary,
  title={\textsc{Exercice}~\thetcbcounter\ifx\relax#1\relax\relax\else\textmd{ --- \itshape #1}\fi}}
\newtcolorbox[auto counter]{corrige}[1][]{boxbase,colback=excol!4,colframe=excol,
  title={\textsc{Corrigé}~\thetcbcounter\ifx\relax#1\relax\relax\else\textmd{ --- \itshape #1}\fi}}
\newcommand{\R}{\mathbb{R}}\newcommand{\N}{\mathbb{N}}\newcommand{\Z}{\mathbb{Z}}\newcommand{\ds}{\displaystyle}
% (un \usepackage{fancyhdr}+en-tête est possible ; il est ignoré côté web)
% =========================================================================
\begin{document}

\begin{corrige}[les exceptions : peroxydes et hydrures]
Dans un \emph{peroxyde} l'oxygène est $-\mathrm{I}$ ; dans un \emph{hydrure} l'hydrogène est $-\mathrm{I}$.
On le retrouve par l'équation de fermeture.
\begin{enumerate}[label=\alph*)]
  \item \ce{H2O2} : $2(+1)+2x=0 \Rightarrow 2x=-2 \Rightarrow x=-\mathrm{I}$ (peroxyde).
  \item \ce{Na2O2} : $2(+1)+2x=0 \Rightarrow x=-\mathrm{I}$ (peroxyde de sodium) — et non $-\mathrm{II}$.
  \item \ce{NaH} : $(+1)_{\ce{Na}}+x=0 \Rightarrow x=-\mathrm{I}$ (hydrure : l'hydrogène joue le rôle de \ce{H-}).
  \item \ce{CaH2} : $(+2)+2x=0 \Rightarrow 2x=-2 \Rightarrow x=-\mathrm{I}$ (hydrure).
\end{enumerate}
\end{corrige}

\begin{corrige}[étage d'oxydation et décompte des électrons]
\begin{enumerate}[label=\alph*)]
  \item \ce{Cu2O} est neutre, $\ce{O}=-\mathrm{II}$ : $2x+(-2)=0 \Rightarrow 2x=+2 \Rightarrow x=+\mathrm{I}$.
        Le cuivre est au degré $+\mathrm{I}$.
  \item Le cuivre neutre a $Z=29$ électrons ; au degré $+\mathrm{I}$ il a \emph{cédé} un électron :
        $29-1=\boxed{28}$ électrons.
  \item L'oxygène ($Z=8$) au degré $-\mathrm{II}$ a \emph{gagné} deux électrons : $8-(-2)=8+2=\boxed{10}$ électrons.
\end{enumerate}
\emph{Rappel du piège :} le cuivre étant un métal, son n.o.\ ne peut pas être négatif — on écarte « $-\mathrm{I}$ ».
\end{corrige}

\begin{corrige}[nombre d'électrons d'un ion monoatomique]
On applique $N_{\ce{e-}} = Z - q$.
\begin{enumerate}[label=\alph*)]
  \item \ce{Al^3+} : $13-(+3)=10$ électrons.
  \item \ce{S^2-} : $16-(-2)=18$ électrons.
  \item \ce{Fe^3+} : $26-(+3)=23$ électrons.
  \item \ce{N^3-} : $7-(-3)=10$ électrons.
\end{enumerate}
\emph{Remarque :} \ce{Al^3+}, \ce{N^3-} et \ce{O^2-} ont tous $10$ électrons (configuration du néon).
\end{corrige}

\begin{corrige}[raisonnement inverse : du n.o.\ à la charge de l'oxoanion]
On écrit toujours : (n.o.\ central) $+$ (nombre de \ce{O})$\times(-2) = $ charge.
\begin{enumerate}[label=\alph*)]
  \item \ce{SO3} avec $\ce{S}=+\mathrm{IV}$ : $(+4)+3(-2)=-2$. La charge est $-2$ : c'est \ce{SO3^2-} (sulfite).
  \item Un oxygène, $\ce{Cl}=+\mathrm{I}$ : $(+1)+(-2)=-1$. La charge est $-1$ : c'est \ce{ClO-} (hypochlorite).
  \item Charge $-1$, $\ce{N}=+\mathrm{V}$ : $(+5)+x(-2)=-1 \Rightarrow -2x=-6 \Rightarrow x=3$.
        L'ion contient $3$ oxygènes : c'est \ce{NO3-} (nitrate).
\end{enumerate}
\end{corrige}

\begin{corrige}[synthèse : les deux oxydes du fer]
\begin{enumerate}[label=\alph*)]
  \item \ce{FeO} : $x+(-2)=0 \Rightarrow x=+\mathrm{II}$. \quad \ce{Fe2O3} : $2x+3(-2)=0 \Rightarrow 2x=+6 \Rightarrow x=+\mathrm{III}$.
  \item \ce{FeO} contient du \textbf{fer(II)} (oxyde de fer(II), ferreux) ; \ce{Fe2O3} contient du
        \textbf{fer(III)} (oxyde de fer(III), ferrique).
  \item Dans \ce{Fe2O3}, le fer est l'ion \ce{Fe^3+} : $N_{\ce{e-}}=26-3=23$ électrons.
\end{enumerate}
\end{corrige}

\end{document}
